\documentclass{article}
\usepackage{amsmath,amssymb,amsthm}
\usepackage{algorithm,algorithmic}
\usepackage{hyperref}
\usepackage{enumitem}
\newtheorem{theorem}{Theorem}
\newtheorem{lemma}{Lemma}
\newtheorem{assumption}{Assumption}
\newcommand{\norm}[1]{\left\lVert#1\right\rVert}
\newcommand{\inner}[1]{\left\langle#1\right\rangle}
\begin{document}

\section*{Heavy‑Ball Gradient Method (HB)}

\section*{Mathematical Formulation}
Let $f:\mathbb{R}^d\to\mathbb{R}$ be differentiable.  
Given step size $\alpha>0$ and momentum parameter $\beta\in[0,1)$, the Heavy‑Ball method generates $\{x_k\}_{k\ge 0}$ by  
\[
\boxed{%
x_{k+1}=x_k+\beta\bigl(x_k-x_{k-1}\bigr)-\alpha\nabla f(x_k),\qquad k\ge 1,
}
\]
with an initial pair $x_0,x_{-1}\in\mathbb{R}^d$ (typically $x_{-1}=x_0$).

\section*{Pseudocode}
\begin{algorithm}[H]
\caption{Heavy‑Ball Gradient Method}
\begin{algorithmic}[1]
\REQUIRE Objective $f$, step size $\alpha>0$, momentum $\beta\in[0,1)$,
initial point $x_0\in\mathbb{R}^d$, optional ``previous'' point $x_{-1}=x_0$,
maximum iterations $K$.
\FOR{$k=0$ \TO $K-1$}
    \STATE Compute gradient $g_k\gets \nabla f(x_k)$
    \STATE Update $x_{k+1}\gets x_k+\beta (x_k-x_{k-1})-\alpha g_k$
\ENDFOR
\RETURN $x_{K}$
\end{algorithmic}
\end{algorithm}

\section*{Dry Run Example}
Consider the one‑dimensional quadratic
\[
f(w)=(w-3)^2,\qquad \nabla f(w)=2(w-3).
\]
Choose $\alpha=0.1$, $\beta=0.5$, and initialise $x_0=0$, $x_{-1}=0$.

\begin{center}
\begin{tabular}{c|c|c|c}
$k$ & $x_k$ & $\nabla f(x_k)$ & $f(x_k)$\\ \hline
0 & $0$ & $2(0-3)=-6$ & $9$\\
1 & $x_1=0+0.5(0-0)-0.1(-6)=0.6$ & $2(0.6-3)=-4.8$ & $(0.6-3)^2=5.76$\\
2 & $x_2=0.6+0.5(0.6-0)-0.1(-4.8)=0.6+0.3+0.48=1.38$ &
$2(1.38-3)=-3.24$ & $(1.38-3)^2\approx2.6244$\\
3 & $x_3=1.38+0.5(1.38-0.6)-0.1(-3.24)=1.38+0.39+0.324=2.094$ &
$2(2.094-3)=-1.812$ & $(2.094-3)^2\approx0.820$\\
\end{tabular}
\end{center}
The objective value decreases from $9$ to $\approx0.82$ after three iterations.

\newpage

\section*{Assumptions}
\begin{assumption}[Smoothness]\label{as:Lsmooth}
$f$ has $L$‑Lipschitz continuous gradient, i.e.
\[
\forall x,y\in\mathbb{R}^d:\quad
\norm{\nabla f(x)-\nabla f(y)}\le L\norm{x-y}.
\]
\end{assumption}
\begin{assumption}[Lower Boundedness]\label{as:lowerbound}
$f$ is bounded below: $f_*:=\inf_{x}f(x)>-\infty$.
\end{assumption}
\begin{assumption}[Parameter Choice]\label{as:params}
The step size $\alpha$ and momentum $\beta$ satisfy
\[
0<\alpha\le \frac{1-\beta}{L},\qquad 0\le\beta<1.
\]
\end{assumption}

\section*{Main Convergence Theorem}
\begin{theorem}\label{thm:hb}
Let $\{x_k\}$ be generated by the Heavy‑Ball method under Assumptions \ref{as:Lsmooth}–\ref{as:params}.  
Define the Lyapunov sequence
\[
\Phi_k\;:=\;f(x_k)-f_*+\frac{\beta}{2\alpha}\,\norm{x_k-x_{k-1}}^2,\qquad k\ge 1.
\]
Then for every $K\ge 1$,
\[
\sum_{k=0}^{K-1}\norm{\nabla f(x_k)}^2
\;\le\;
\frac{2}{\alpha(1-\beta)}\bigl(\Phi_0-\Phi_K\bigr)
\;\le\;
\frac{2}{\alpha(1-\beta)}\bigl(f(x_0)-f_*\bigr).
\]
Consequently,
\[
\min_{0\le k\le K-1}\norm{\nabla f(x_k)}^2
\;\le\;
\frac{2}{\alpha(1-\beta)K}\bigl(f(x_0)-f_*\bigr)
\;=\;O\!\left(\frac{1}{K}\right).
\]
Thus the algorithm attains the standard non‑convex $O(1/K)$ rate for the squared gradient norm.
\end{theorem}

\section*{Theoretical Properties}
\begin{itemize}[leftmargin=*]
\item \textbf{Stability.} The Lyapunov function $\Phi_k$ is non‑increasing because the descent inequality in Theorem \ref{thm:hb} holds for any $k$. Hence the iterates remain in a level set of $f$ and are bounded.
\item \textbf{Hyper‑parameter sensitivity.} Condition \eqref{as:params} forces $\alpha$ to shrink as $\beta$ grows; the product $\alpha(1-\beta)$ appears in the denominator of the rate, showing a trade‑off: larger momentum accelerates progress only if the step size is reduced accordingly.
\item \textbf{Comparison to baselines.}
\begin{enumerate}
\item With $\beta=0$ the method reduces to plain gradient descent and the bound becomes $\frac{2}{\alpha}\bigl(f(x_0)-f_*\bigr)/K$, the classic $O(1/K)$ result.
\item For $0<\beta<1$, the same $O(1/K)$ order is retained, but the constant improves by a factor $1/(1-\beta)$, reflecting the known acceleration effect of the heavy‑ball momentum in practice.
\end{enumerate}
\end{itemize}

\section*{Discussion}
The theorem shows that, despite the lack of convexity, the Heavy‑Ball scheme enjoys the same worst‑case convergence order as gradient descent when the step size respects the $L$‑smoothness and the momentum parameter. The proof hinges on a carefully chosen Lyapunov function that couples function values with the kinetic energy $\norm{x_k-x_{k-1}}^2$. The result also clarifies why overly aggressive momentum (large $\beta$) without a commensurate reduction of $\alpha$ can violate Assumption \ref{as:params} and lead to divergence.

\newpage

\section*{Preliminaries (Definitions and Lemmas)}
\begin{lemma}[Smoothness Inequality]\label{lem:smooth}
If $f$ satisfies Assumption \ref{as:Lsmooth}, then for any $x,y$,
\[
f(y)\le f(x)+\inner{\nabla f(x)}{y-x}+\frac{L}{2}\norm{y-x}^2.
\]
\end{lemma}
\begin{proof}
Standard consequence of $L$‑Lipschitz gradient; see e.g. Nesterov (2004). \hfill $\square$
\end{proof}

\section*{Main Proof (Step‑by‑Step Derivation)}
We prove Theorem \ref{thm:hb}.  All algebraic steps are displayed and each non‑trivial manipulation is labelled.

\bigskip
\noindent\textbf{Notation.} For brevity write $g_k:=\nabla f(x_k)$ and $s_k:=x_k-x_{k-1}$.

\noindent\textbf{Step 1 (Expand the update).}  
From the algorithm,
\[
x_{k+1}=x_k+\beta s_k-\alpha g_k. \tag{1}
\]

\noindent\textbf{Step 2 (Smoothness bound on $f(x_{k+1})$).}  
Apply Lemma \ref{lem:smooth} with $x=x_k$, $y=x_{k+1}$:
\[
f(x_{k+1})\le f(x_k)+\inner{g_k}{x_{k+1}-x_k}+\frac{L}{2}\norm{x_{k+1}-x_k}^2.
\tag{2}
\]

\noindent\textbf{Step 3 (Insert the update (1) into (2)).}  
Since $x_{k+1}-x_k=\beta s_k-\alpha g_k$,
\[
\begin{aligned}
f(x_{k+1}) &\le f(x_k)+\inner{g_k}{\beta s_k-\alpha g_k}
            +\frac{L}{2}\norm{\beta s_k-\alpha g_k}^2\\
&= f(x_k)+\beta\inner{g_k}{s_k}-\alpha\norm{g_k}^2
   +\frac{L}{2}\bigl(\beta^2\norm{s_k}^2-2\alpha\beta\inner{g_k}{s_k}+\alpha^2\norm{g_k}^2\bigr).
\end{aligned}
\tag{3}
\]

\noindent\textbf{Step 4 (Collect like terms).}  
Group the coefficients of $\inner{g_k}{s_k}$ and $\norm{g_k}^2$:
\[
\begin{aligned}
f(x_{k+1}) &\le f(x_k)
   +\bigl(\beta-\alpha L\beta\bigr)\inner{g_k}{s_k}
   -\bigl(\alpha-\tfrac{L\alpha^2}{2}\bigr)\norm{g_k}^2
   +\tfrac{L\beta^2}{2}\norm{s_k}^2 .
\end{aligned}
\tag{4}
\]

\noindent\textbf{Step 5 (Use the parameter condition).}  
Assumption \ref{as:params} gives $\alpha\le\frac{1-\beta}{L}$, which implies
\[
\alpha-\frac{L\alpha^2}{2}\;\ge\;\frac{\alpha}{2}(1-\beta). \tag{5}
\]
Also $0\le\beta<1$ yields $0\le\beta-\alpha L\beta\le\beta$; we keep the exact term for now.

\noindent\textbf{Step 6 (Introduce the kinetic term).}  
Add and subtract $\dfrac{\beta}{2\alpha}\norm{s_{k+1}}^2$ to the right‑hand side, where $s_{k+1}=x_{k+1}-x_k=\beta s_k-\alpha g_k$:
\[
\frac{\beta}{2\alpha}\norm{s_{k+1}}^2
   =\frac{\beta}{2\alpha}\bigl(\beta^2\norm{s_k}^2-2\alpha\beta\inner{g_k}{s_k}+\alpha^2\norm{g_k}^2\bigr).
\tag{6}
\]

\noindent\textbf{Step 7 (Form the Lyapunov difference).}  
Define $\Phi_k:=f(x_k)-f_*+\dfrac{\beta}{2\alpha}\norm{s_k}^2$.  Using (4)–(6) we obtain
\[
\begin{aligned}
\Phi_{k+1}
&=f(x_{k+1})-f_*+\frac{\beta}{2\alpha}\norm{s_{k+1}}^2\\
&\le f(x_k)-f_*+\frac{\beta}{2\alpha}\norm{s_k}^2
   +\bigl(\beta-\alpha L\beta\bigr)\inner{g_k}{s_k}
   -\bigl(\alpha-\tfrac{L\alpha^2}{2}\bigr)\norm{g_k}^2
   +\tfrac{L\beta^2}{2}\norm{s_k}^2 \\
&\qquad\;-\frac{\beta}{2\alpha}\bigl(\beta^2\norm{s_k}^2-2\alpha\beta\inner{g_k}{s_k}+\alpha^2\norm{g_k}^2\bigr)\\[4pt]
&= \Phi_k
   -\underbrace{\Bigl(\alpha-\tfrac{L\alpha^2}{2}-\tfrac{\beta\alpha}{2}\Bigr)}_{\displaystyle\ge\frac{\alpha}{2}(1-\beta)}
      \norm{g_k}^2
   +\underbrace{\Bigl(\tfrac{L\beta^2}{2}-\tfrac{\beta^3}{2\alpha}\Bigr)}_{\displaystyle\le 0}
      \norm{s_k}^2 .
\end{aligned}
\tag{7}
\]

\noindent\textbf{Step 8 (Sign of the $s_k$ term).}  
From $\alpha\le\frac{1-\beta}{L}$ we have $\frac{L\beta^2}{2}\le\frac{\beta^2(1-\beta)}{2\alpha}\le\frac{\beta^3}{2\alpha}$, hence the coefficient of $\norm{s_k}^2$ in (7) is non‑positive. Dropping this non‑positive term yields a clean descent inequality:
\[
\Phi_{k+1}\le \Phi_k-\frac{\alpha}{2}(1-\beta)\norm{g_k}^2.
\tag{8}
\]

\noindent\textbf{Step 9 (Summation).}  
Summing (8) from $k=0$ to $K-1$ gives
\[
\Phi_K\le \Phi_0-\frac{\alpha}{2}(1-\beta)\sum_{k=0}^{K-1}\norm{g_k}^2.
\tag{9}
\]
Since $\Phi_K\ge 0$ (by definition $f\ge f_*$ and the kinetic term is non‑negative), we rearrange (9):
\[
\sum_{k=0}^{K-1}\norm{g_k}^2\le \frac{2}{\alpha(1-\beta)}\Phi_0
   \le \frac{2}{\alpha(1-\beta)}\bigl(f(x_0)-f_*\bigr).
\tag{10}
\]

\noindent\textbf{Step 10 (Deriving the rate).}  
Divide (10) by $K$ and use $\min_{0\le k\le K-1}\norm{g_k}^2\le \frac{1}{K}\sum_{k=0}^{K-1}\norm{g_k}^2$:
\[
\min_{0\le k\le K-1}\norm{\nabla f(x_k)}^2
\;\le\;
\frac{2}{\alpha(1-\beta)K}\bigl(f(x_0)-f_*\bigr).
\tag{11}
\]
Equation (11) is exactly the $O(1/K)$ bound claimed in Theorem \ref{thm:hb}.  $\square$

\section*{Rate Derivation (With Constants)}
The constant hidden in the $O(\cdot)$ notation is explicitly
\[
C:=\frac{2\bigl(f(x_0)-f_*\bigr)}{\alpha(1-\beta)}.
\]
Thus for any $K\ge 1$,
\[
\min_{0\le k\le K-1}\norm{\nabla f(x_k)}^2\le \frac{C}{K}.
\]

\section*{Special Cases}
\begin{enumerate}
\item \textbf{Pure Gradient Descent ($\beta=0$).}  
The Lyapunov reduces to $\Phi_k=f(x_k)-f_*$ and (8) becomes
$f(x_{k+1})\le f(x_k)-\frac{\alpha}{2}\norm{\nabla f(x_k)}^2$, the classical descent lemma.

\item \textbf{Quadratic Objective.}  
If $f(x)=\frac12 x^\top Q x$ with $Q\succeq 0$, the Heavy‑Ball dynamics are linear and the condition $\alpha\le\frac{1-\beta}{\lambda_{\max}(Q)}$ is both necessary and sufficient for linear convergence. The present proof still holds and yields the same $O(1/k)$ bound, which is sub‑optimal for quadratics (the true rate is geometric).

\item \textbf{Large Momentum Limit.}  
Let $\beta\to 1^{-}$ while keeping $\alpha(1-\beta)=\tilde\alpha$ constant. Then the bound (11) becomes
$\min\|\nabla f\|^2\le \frac{2}{\tilde\alpha K}(f(x_0)-f_*)$, showing that the effective step size is $\tilde\alpha$, confirming the well‑known “effective stepsize = $\alpha/(1-\beta)$’’ intuition.
\end{enumerate}

\end{document}